\section{Preliminaries}

A cellular automaton is a collection of cells arranged in a lattice, where each cell takes a state from a finite set. The system evolves in discrete time steps and, at each step, all cells update their state synchronously as a function of the states of their neighbors~\cite{weisstein2002_ca,wolfram1982_ca,fsancho2016_ca}.

\begin{definition}
\label{def:ca}
A cellular automaton is a 4-tuple $(\mathcal{L}, \mathcal{S}, \mathcal{N}, \mathcal{R})$ where:
\begin{itemize}
    \item $\mathcal{L} \subseteq \mathbb{Z}^D$ is the $D$-dimensional cellular space; each cell is identified by its position $\vec{i} \in \mathcal{L}$.

    \item $\mathcal{S} = \{0, 1, \dots, q-1\}$ is the finite set of $q$ states. The state of cell $\vec{i}$ at time $t$ is denoted by $\sigma_{\vec{i}}(t) \in \mathcal{S}$.

    \item $\mathcal{N} = (\vec{v}_1, \dots, \vec{v}_m)$ is the ordered tuple of the $m$ displacement vectors that define the neighborhood: the neighbors of $\vec{i}$ are the cells $\vec{i} + \vec{v}_j$, with $j = 1, \dots, m$.

    \item $\mathcal{R} : \mathcal{S}^m \to \mathcal{S}$ is the local rule, which assigns the new state of a cell from the states of its $m$ neighbors:
    \[
    \sigma_{\vec{i}}(t + 1) = \mathcal{R}\left( \sigma_{\vec{i} + \vec{v}_1}(t), \dots, \sigma_{\vec{i} + \vec{v}_m}(t) \right).
    \]
\end{itemize}
\end{definition}

The order of the tuple $\mathcal{N}$ matters: it determines the numbering of the rules and allows us to define relations such as mirror rules (Section~\ref{subsection:rules_same_domain}).

\subsection{Space, neighborhoods and rules}

In practice, the cellular space is restricted by open boundaries, which assign a fixed state to the missing neighbors, or periodic ones, which take them from the opposite boundary~\cite{bhattacharjee2019_ca,ponce2015_ca}. Throughout this work we always use periodic boundaries: in one dimension, the space is a ring of $n$ cells, $\mathcal{L} = \mathbb{Z}_n$, and with radius $r$ the neighborhood contains the $m = 2r + 1$ cells at distance at most $r$, including the cell itself.

The elementary cellular automaton (ECA) has $q = 2$, radius $r = 1$ ($m = 3$) and rule $\mathcal{R} : \{0,1\}^3 \to \{0,1\}$. There are $q^{m} = 2^3 = 8$ possible neighborhood patterns and each one is assigned a new state, so the total number of rules is $q^{q^{m}} = 2^8 = 256$. Each rule is identified by a decimal number obtained by reading, as a binary number, the states assigned to the patterns $111, 110, \dots, 000$~\cite{Wolfram2002_nks}. In general, we denote by $r_{q,m}$ the rule number $r$ with $q$ states and $m$ neighbors, numbered in the same way in base $q$; thus, the ECA are the rules $r_{2,3}$. When $m$ is even the neighborhood cannot be symmetric: with $m = 2$ we use the cells $(i-1, i)$, which is the \textit{Mathematica} convention for radius $r = 1/2$.

\subsection{Global function and transition graph}

Let $\mathcal{C} = \mathcal{S}^{\mathcal{L}}$ be the set of configurations, with $N = |\mathcal{C}| = q^{|\mathcal{L}|}$. The global function $F : \mathcal{C} \to \mathcal{C}$ applies the local rule to all cells synchronously. The transition graph $G_F = (\mathcal{C}, E)$ has one node per configuration and a directed edge $\sigma \to F(\sigma)$. Since $F$ is a function, every node has out-degree 1 (a functional graph). Each connected component contains exactly one cycle, the attractor, with transient trees hanging from it; together they form its basin of attraction. For each $x \in \mathcal{C}$ we define its period $p(x)$ as the length of the cycle it reaches, and its height $h(x)$ as the minimum number of steps needed to reach it ($h(x) = 0$ if $x$ is periodic). Configurations without a predecessor are known as Garden of Eden configurations.

\subsection{Quotient by translations}
\label{subsec:quotient}

Let $T$ be the translation by one position in $\mathcal{L}$ (a rotation of the ring $\mathbb{Z}_n$). With periodic boundaries, $F \circ T = T \circ F$, so two configurations that differ by a translation have equivalent trajectories. This allows us to build the quotient graph, whose nodes are the classes of configurations under translation and whose edges are $[x] \to [F(x)]$. All configurations in the same class have the same height, and the number of nodes drops to approximately $N/|\mathcal{L}|$. In the quotient graph each class counts as a single node, so the proportions computed on it are normalized by the number of classes.
